%% ma: L11-B09-0002
%% lop: 11
%% chuong: 3
%% bai: 9
%% dang: 11.9.2
%% loai: TN4
%% muc: TH
%% dap_an: "B"
%% nguon: "(NBV)-ĐỀ SỐ 2-MỖI NGÀY 1 ĐỀ THI 2025.pdf (D02, MỖI NGÀY 1 ĐỀ - Đề số 2), Phần I Câu 12"
%% kien_thuc: "Bài 9 (mốt của mẫu số liệu ghép nhóm)"
%% trang_thai: cho-duyet
%% ly_do: "Dạng toán mới đề xuất 11.9.2 (mốt của mẫu số liệu ghép nhóm, so sánh hai mẫu), chờ giáo viên duyệt."
\begin{ex}
Cho hai mẫu số liệu ghép nhóm $A$ và $B$ có bảng tần số ghép nhóm như sau:
\begin{center}
\begin{tabular}{|c|c|c|c|c|c|}
\hline
$A$ & $[160;180)$ & $[180;200)$ & $[200;220)$ & $[220;240)$ & $[240;260)$\\
\hline
Tần số & $14$ & $22$ & $8$ & $6$ & $4$\\
\hline
\end{tabular}

\medskip
\begin{tabular}{|c|c|c|c|c|c|}
\hline
$B$ & $[16;18)$ & $[18;20)$ & $[20;22)$ & $[22;24)$ & $[24;26)$\\
\hline
Tần số & $7$ & $11$ & $4$ & $3$ & $2$\\
\hline
\end{tabular}
\end{center}
Gọi $M_0^A$ và $M_0^B$ lần lượt là mốt của mẫu số liệu ghép nhóm $A$ và $B$. Phát biểu nào sau đây là đúng?
\choice{$M_0^A=2M_0^B$}{\True $M_0^A=10M_0^B$}{$M_0^A=20M_0^B$}{$M_0^A=M_0^B+144$}
\loigiai{Nhóm chứa mốt của $A$ là $[180;200)$:
$M_0^A=180+\dfrac{22-14}{(22-14)+(22-8)}\cdot20=\dfrac{2060}{11}$.
Nhóm chứa mốt của $B$ là $[18;20)$:
$M_0^B=18+\dfrac{11-7}{(11-7)+(11-4)}\cdot2=\dfrac{206}{11}$.
Vậy $M_0^A=10M_0^B$.}
\end{ex}
